\subsubsection*{Chapter 1 - Sheaves}

\textcolor{cyan}{Atiyah-Macdonald Exercise 2.15} is used thoroughly. \textcolor{cyan}{(i)} means $t\in\mathcal{F}_x$ can be represented by $s\in\mathcal{F}(U)$, \textcolor{cyan}{(ii)} is used for finding agreement on a smaller neighborhood.

The restriction $\mathcal{F}|_Z$ can be simplified if $Z$ is an open set.

\subsubsection*{Chapter 2 - Schemes}

Most exercises help to enhance the comprehension of the definitions.

\subsubsection*{Chapter 3 - First Properties of Schemes}

The definitions generalize the results in commutative algebra. Gluing morphisms and schemes is a new technique and is used through the exercises.

There is an important theorem which is proved in \Cref{ex3.8}: The intersection of two affine subschemes can be covered by $\left\{W_i\right\}$, where each $W_i$ is basic open in both subschemes.

\Cref{ex3.11.2} is useful when deal with closed subschemes.

\subsubsection*{Chapter 4 - Separated and Proper Morphisms}

\Cref{4.4} and \Cref{4.8} are fundamental for understanding primary ideal structure.

\Cref{4.9} is useful for computations.

\subsubsection*{Chapter 5 - Integral Dependence and Valuations}

The most important chapter in the book.

\textbf{Important results:} \hyperref[5.10]{Lying Over Theorem}, \hyperref[noethernormalization]{Noether Normalization Lemma}.

\Cref{ex5.10} connects the going-up and going-down properties of a ring homomorphism with topological properties of the induced map on spectra, but the ideas are not used further in the book.

\Cref{ex5.25} uses Jacobson rings to generalize Zariski's Lemma.

\subsubsection*{Chapter 6 \& 7 - Chain Conditions \& Noetherian Rings}

Most exercises further refine old concepts rather than introducing new ideas.

\subsubsection*{Chapter 8 - Artinian Rings}

An important result is \Cref{8.5}.

\Cref{ex8.3} provides another generalization of Zariski's Lemma.

\subsubsection*{Chapter 9 - Discrete Valuation Rings and Dedekind Domains}

This chapter is closely connected to algebraic number theory. \Cref{9.2} and \Cref{9.3} appear repeatedly throughout the exercises.
